% =====================================================================
% Appendix H — Physicochemical literature data
% =====================================================================
\chapter{Physicochemical Literature Datasets}
\label{chap:app-physchem}

Sources: \code{solubility\_literature.csv} (16 determinations) and \code{permeability\_literature.csv} (9 records), reproduced in full; these underpin the BCS classification of \cref{sec:lit-classification}.

\section{Solubility}
\small\setlength{\tabcolsep}{4pt}
\begin{longtable}{p{2.4cm}p{2.2cm}p{2.6cm}p{1.6cm}p{3.4cm}}
\caption{Solubility literature dataset.}\label{tab:solubility-full}\\
\toprule
Solvent / medium & Form & Solubility & $T$ & Reference (notes)\\
\midrule
\endfirsthead
\toprule
Solvent / medium & Form & Solubility & $T$ & Reference\\
\midrule
\endhead
\midrule
\multicolumn{5}{r}{\emph{continued on next page}}\\
\endfoot
\bottomrule
\endlastfoot
Water & Tobramycin & $\geq$100 mg/mL ($\geq$213.9 mM) & RT & BenchChem 2026\\
Water & Tobramycin & 94 mg/mL (201.1 mM) & 25\,\textdegree C & \citet{Chen2009}\\
Water & Tobramycin & 50 mg/mL & RT & Sigma-Aldrich\\
Water & Tobr.\ sulfate & 100 mg/mL & RT & USP monograph\\
Water & Tobr.\ sulfate & 50 mg/mL & RT & Various\\
PBS pH 7.2 & Tobramycin & 10 mg/mL & RT & PBPK study\\
DMSO & Tobramycin & 2 mg/mL & RT & \citet{Chen2009} (conflicting)\\
DMSO & Tobramycin & Insoluble & RT & Multiple (contradicts above)\\
Ethanol & Tobramycin & Insoluble & RT & USP\\
Ethanol & Tobr.\ sulfate & Slightly soluble & RT & USP\\
Methanol & Tobr.\ sulfate & Slightly soluble & RT & USP\\
Acetone & Tobr.\ sulfate & Slightly soluble & RT & USP\\
FaSSIF & Tobramycin & Dissolved & 37\,\textdegree C & Predicted from solubility\\
FeSSIF & Tobramycin & Dissolved & 37\,\textdegree C & Predicted from solubility\\
SGF pH 1.2 & Tobramycin & High & 37\,\textdegree C & Predicted (BCS)\\
SIF pH 6.8 & Tobramycin & High & 37\,\textdegree C & Predicted (BCS)\\
\end{longtable}

\section{Permeability}
\small\setlength{\tabcolsep}{4pt}
\begin{longtable}{p{2.8cm}p{2.4cm}p{2.6cm}p{1.6cm}p{3.6cm}}
\caption{Permeability literature dataset (reference year standardized to 2023; the source file contains a typographic ``203'').}\label{tab:perm-full}\\
\toprule
Assay / endpoint & Parameter & Value & Units & Reference (notes)\\
\midrule
\endfirsthead
\toprule
Assay / endpoint & Parameter & Value & Units & Reference\\
\midrule
\endhead
\midrule
\multicolumn{5}{r}{\emph{continued on next page}}\\
\endfoot
\bottomrule
\endlastfoot
Caco-2 & $P_{\mathrm{app}}$ A-to-B & Low, $<1\times10^{-6}$ & cm/s & \citet{Asad2023}; class III confirmation\\
Caco-2 & $P_{\mathrm{app}}$ B-to-A & n.d. & --- & Efflux ratio not determined; polycationic substrate\\
MDCK-MDR1 & Permeability & Low (estimated) & --- & From BCS class, literature\\
PAMPA & Permeability & Low & --- & pH 6.5, literature\\
Ussing chamber & Permeability & Low & --- & Jejunum tissue, literature\\
Oral bioavailability & Absolute $\F$ & 1--2 & \% & IV reference, literature estimate\\
Efflux (P-gp) & Substrate? & Conflicting/minimal & --- & Minimal interaction\\
Efflux (BCRP) & Substrate? & Unknown & --- & ---\\
Influx (OCT2) & Substrate & Possibly & --- & Renal tubular handling\\
\end{longtable}

\noindent\textbf{Reconciliation note.} Two records in the source files carry inconsistencies, documented here for transparency: (i) the DMSO solubility conflict (2~mg/mL vs.\ insoluble) is immaterial to BCS classification, which rests on the aqueous/biorelevant data; (ii) oral bioavailability appears as ``$<$5\%'' (generic envelope) and ``1--2\%'' (point estimate) --- reconciled in \cref{sec:lit-pk}.
